\documentclass[11pt,a4paper]{article}

\usepackage[T1]{fontenc}
\usepackage[utf8]{inputenc}
\usepackage{newtxtext,newtxmath}
\usepackage[a4paper,left=1.7cm,right=1.7cm,top=1.45cm,bottom=1.55cm]{geometry}
\usepackage{microtype}
\usepackage{amsmath}
\usepackage{xcolor}
\usepackage{tabularx}
\usepackage{array}
\usepackage{enumitem}
\usepackage[hidelinks]{hyperref}
\usepackage{fancyhdr}

\definecolor{sectiongray}{gray}{0.42}
\definecolor{muted}{gray}{0.28}

\setlength{\parindent}{0pt}
\setlength{\parskip}{0pt}
\renewcommand{\arraystretch}{1.02}
\pagestyle{fancy}
\fancyhf{}
\renewcommand{\headrulewidth}{0pt}

\newcommand{\cvsection}[1]{%
  \vspace{0.82em}%
  \noindent
  \begin{tabularx}{\textwidth}{@{}>{\raggedleft\arraybackslash}p{2.55cm}@{\hspace{0.35cm}}X@{}}
    \textcolor{sectiongray}{\rule{2.35cm}{1.15pt}} & {\large\color{sectiongray}#1}
  \end{tabularx}
  \vspace{0.18em}%
}

\newcommand{\datedentry}[2]{%
  \noindent
  \begin{tabularx}{\textwidth}{@{}>{\raggedleft\arraybackslash}p{2.55cm}@{\hspace{0.35cm}}X@{}}
    \small #1 & \parbox[t]{\linewidth}{#2}
  \end{tabularx}
  \vspace{0.15em}%
}

\newenvironment{preprints}{%
  \begin{enumerate}[label=\arabic*.,leftmargin=3.28cm,labelsep=0.55em,itemsep=0.42em,topsep=0.15em,parsep=0pt]
}{%
  \end{enumerate}
}

\begin{document}

\begin{tabularx}{\textwidth}{@{}X>{\raggedleft\arraybackslash}p{7.0cm}@{}}
  {\fontsize{25}{29}\selectfont Leonardo Franchi} &
  \begin{tabular}{@{}r@{}}
    PhD Student in Mathematics\\
    University of Cambridge\\
    \href{mailto:lf511@cam.ac.uk}{lf511@cam.ac.uk}\\
    \href{https://leo-franchi.github.io/}{leo-franchi.github.io}
  \end{tabular}
\end{tabularx}

\vspace{-0.15em}
{\large\itshape Curriculum Vitae}

\cvsection{Education}

\datedentry{Oct. 2025--present}{%
  \textbf{PhD in Mathematics}, University of Cambridge, Cambridge, UK.\\
  Supervisor: Timothy Gowers.
}

\datedentry{Oct. 2023--Jul. 2025}{%
  \textbf{Master's Degree in Mathematics}, University of Padua, Padua, Italy.\\
  Final grade: 110/110 cum laude. Thesis: \emph{Osgood meets DiPerna--Lions}. Advisor: Guido De Philippis.
}

\datedentry{Oct. 2020--Oct. 2025}{%
  \textbf{Galilean Diploma}, Galilean School of Higher Education, University of Padua, Padua, Italy.
}

\datedentry{Oct. 2020--Feb. 2023}{%
  \textbf{Bachelor's Degree in Mathematics}, University of Padua, Padua, Italy.\\
  Final grade: 110/110 cum laude.
}

\cvsection{Submitted papers and preprints}

\begin{preprints}
  \item Guido De Philippis and Leonardo Franchi,
  \emph{Osgood meets Ambrosio--DiPerna--Lions},
  \href{https://arxiv.org/abs/2607.28118}{arXiv:2607.28118} (2026).

  \item Valentino Badalucco and Leonardo Franchi,
  \emph{Fractal uncertainty principle over $\mathbb{Q}_p$},
  \href{https://arxiv.org/abs/2607.11534}{arXiv:2607.11534} (2026).

  \item Leonardo Franchi, W. T. Gowers and Fredy Yip,
  \emph{Product-free subsets of $(0,1)$},
  \href{https://arxiv.org/abs/2607.06073}{arXiv:2607.06073} (2026).

  \item Gaia Carenini and Leonardo Franchi,
  \emph{A unified abstract regularity lemma},
  \href{https://arxiv.org/abs/2606.06192}{arXiv:2606.06192} (2026).

  \item Gaia Carenini and Leonardo Franchi,
  \emph{A strengthening of Chang's lemma},
  \href{https://arxiv.org/abs/2605.07916}{arXiv:2605.07916} (2026).
\end{preprints}

\cvsection{Invited Talks}

\datedentry{12 Oct. 2026}{%
  \emph{Product-free subsets of $(0,1)$ and a strengthened Brunn--Minkowski inequality}.\\
  Junior Number Theory Seminar, University of Oxford, Oxford, UK.
}

\datedentry{29 Sep. 2026}{%
  \emph{Product-free subsets of $(0,1)$ and a strengthened Brunn--Minkowski inequality}.\\
  Combinatorics Seminar, University of Bristol, Bristol, UK.
}

\datedentry{26--28 Jun. 2026}{%
  \emph{A strengthening of Chang's lemma}.\\
  WOMBL 18 (Warwick--Oxbridge--Manchester--Bristol--Lancaster Meeting on Additive Combinatorics and Analytic Number Theory), Ambleside Campus, University of Cumbria, UK.
}

\newpage
\cvsection{Awards and Funding}

\datedentry{2025-2029}{Trinity Cambridge Research Studentship, supported by the Isaac Newton Trust.}
\datedentry{2024}{First Prize, International Mathematics Competition for University Students, Blagoevgrad, Bulgaria.}
\datedentry{2020}{First place, Italian Mathematical Olympiad.}

\cvsection{Teaching and Academic Experiences}

\datedentry{Michaelmas 2026}{%
  \textbf{Undergraduate Supervisor, Optimisation}, Queens' College, University of Cambridge.
}

\datedentry{Lent 2026}{%
  \textbf{Undergraduate Supervisor, Probability (Part IA)}, Queens' College, University of Cambridge.
}

\datedentry{Dec. 2025}{%
  \textbf{Examiner for Undergraduate Mathematics Admissions}, Queens' College, University of Cambridge.
}

\datedentry{Jun. 2023, 2024}{%
  \textbf{Lecturer, International Mathematics Summer Camp}, Beijing Institute of Mathematical Sciences and Applications (BIMSA), Beijing, China.
}

\end{document}
